\section{Definición del problema}
\label{sec:problema}

\subsection{Enunciado del problema central}

\begin{cajaclave}[title={Problema central}]
Las organizaciones que ya poseen sensores IoT no logran convertir los datos que estos generan en
información utilizable, porque carecen de una plataforma que unifique la ingesta, la normalización y la
gestión de dispositivos que se comunican mediante tecnologías y formatos distintos entre sí.
\end{cajaclave}

Formulado desde la perspectiva del afectado, el problema se expresa así: \enquote{tengo los sensores
instalados y funcionando, pero no tengo dónde recibir sus datos, ni cómo administrarlos, ni cómo saber
si siguen reportando}.

\subsection{Análisis de causas}

\begin{longtable}{@{}L{1.1cm}L{3.4cm}L{9.2cm}@{}}
\filaencab \encab{} & \encab{Causa} & \encab{Descripción} \\
\midrule
\endfirsthead
\filaencab \encab{} & \encab{Causa} & \encab{Descripción} \\ \midrule \endhead
\textbf{C1} & Fragmentación de protocolos y formatos &
Cada fabricante define su propio esquema de datos. Un sensor LoRaWAN entrega una trama binaria
propietaria; uno WiFi publica un objeto JSON con nombres de campo arbitrarios; un tercero expone una
interfaz distinta. No existe vocabulario común entre ellos. \\
\textbf{C2} & Acoplamiento entre hardware y plataforma &
La oferta comercial regional vende solución cerrada llave en mano: el panel funciona con el hardware
del propio proveedor. Un sensor ya instalado de otra marca no se integra. \\
\textbf{C3} & Costo y dependencia de plataformas comerciales &
Las alternativas de mercado cobran licencia por dispositivo o por volumen de datos, y alojan la
información en infraestructura de terceros fuera del control de la organización. \\
\textbf{C4} & Ausencia de capacidad técnica interna &
Construir la capa de ingesta, normalización y almacenamiento exige competencias de redes, protocolos y
desarrollo que las organizaciones destinatarias no poseen internamente. \\
\textbf{C5} & Falta de visibilidad del estado de la conectividad &
Sin una capa de gestión no se detecta que un sensor dejó de reportar, que su batería se agotó o que el
enlace de radio se degradó. \\
\bottomrule
\caption{Causas identificadas del problema central.}
\label{tab:causas}
\end{longtable}

\begin{cajaclave}
La causa C5 admite una formulación especialmente compacta, que además orienta el diseño de la
solución: \textbf{el silencio de un sensor es indistinguible de la ausencia de novedades.}
\end{cajaclave}

\subsection{Efectos observables}

\begin{longtable}{@{}L{1.1cm}L{5.4cm}L{7.2cm}@{}}
\filaencab \encab{} & \encab{Efecto} & \encab{Consecuencia} \\ \midrule \endfirsthead
\filaencab \encab{} & \encab{Efecto} & \encab{Consecuencia} \\ \midrule \endhead
\textbf{E1} & Inversión en sensores subutilizada o inmovilizada & Equipamiento adquirido que no produce retorno \\
\textbf{E2} & Datos que se pierden o nunca se almacenan & Imposibilidad de análisis histórico o de tendencias \\
\textbf{E3} & Decisiones operativas tomadas sin dato & El sensor existe, pero no informa la decisión \\
\textbf{E4} & Fallas de sensores detectadas tardíamente & Pérdida de continuidad del monitoreo sin advertencia \\
\textbf{E5} & Dependencia de proveedor & Cada ampliación obliga a comprar al mismo proveedor \\
\bottomrule
\caption{Efectos observables del problema.}
\label{tab:efectos}
\end{longtable}

\begin{figure}[htbp]
\centering
\begin{tikzpicture}[node distance=6mm, font=\scriptsize]

% ---- EFECTOS ----
\node[caja, font=\scriptsize, text width=2.0cm] (e1) {\textbf{E1}\\Inversión\\inmovilizada};
\node[caja, font=\scriptsize, text width=2.0cm, right=3mm of e1] (e2) {\textbf{E2}\\Datos\\perdidos};
\node[caja, font=\scriptsize, text width=2.0cm, right=3mm of e2] (e3) {\textbf{E3}\\Decisiones\\sin dato};
\node[caja, font=\scriptsize, text width=2.0cm, right=3mm of e3] (e4) {\textbf{E4}\\Fallas vistas\\tarde};
\node[caja, font=\scriptsize, text width=2.0cm, right=3mm of e4] (e5) {\textbf{E5}\\Dependencia\\de proveedor};

\begin{scope}[on background layer]
\node[grupo, fit=(e1)(e5)] (gefectos) {};
\end{scope}
\node[etiqueta, left=2mm of gefectos, font=\scriptsize\bfseries\color{gristexto}] {EFECTOS};

% ---- PROBLEMA CENTRAL ----
\node[cajaacento, below=11mm of gefectos, text width=12.0cm, fill=verdeacento!6,
      font=\scriptsize] (prob)
  {\textbf{Quien ya posee sensores heterogéneos no puede convertir sus datos en información
   utilizable:}\\ falta una plataforma que unifique la ingesta, la normalización y la gestión de
   dispositivos.};
\node[etiqueta, left=2mm of prob, font=\scriptsize\bfseries\color{gristexto}] {PROBLEMA\\CENTRAL};

% ---- CAUSAS ----
\node[caja, font=\scriptsize, below=11mm of prob, text width=2.0cm, xshift=-4.65cm] (c1)
  {\textbf{C1}\\Fragmentación\\de protocolos};
\node[caja, font=\scriptsize, text width=2.0cm, right=3mm of c1] (c2) {\textbf{C2}\\Acoplamiento\\hardware};
\node[caja, font=\scriptsize, text width=2.0cm, right=3mm of c2] (c3) {\textbf{C3}\\Costo y\\dependencia};
\node[caja, font=\scriptsize, text width=2.0cm, right=3mm of c3] (c4) {\textbf{C4}\\Sin capacidad\\técnica};
\node[caja, font=\scriptsize, text width=2.0cm, right=3mm of c4] (c5) {\textbf{C5}\\Sin visibilidad\\del enlace};

\begin{scope}[on background layer]
\node[grupo, fit=(c1)(c5)] (gcausas) {};
\end{scope}
\node[etiqueta, left=2mm of gcausas, font=\scriptsize\bfseries\color{gristexto}] {CAUSAS};

% ---- FLECHAS ----
\draw[flechagruesa] (gcausas.north) -- node[nota, right=1mm, midway, font=\scriptsize] {provocan} (prob.south);
\draw[flechagruesa] (prob.north) -- node[nota, right=1mm, midway, font=\scriptsize] {produce} (gefectos.south);

\end{tikzpicture}
\caption{Árbol de problemas: las cinco causas provocan el problema central, que a su vez produce
los cinco efectos observables.}
\label{fig:arbol}
\end{figure}


\subsection{Justificación de relevancia e impacto}

\subsubsection{Relevancia para el propietario del proyecto}

La \datoUniversidad{} dispone de sensores desplegados y requiere un destino para sus datos: una interfaz
de programación a la cual enviarlos y una plataforma desde la cual administrarlos. Esta necesidad,
transmitida por el profesor guía en la reunión del 17 de agosto de 2026, constituye la solicitud que
origina el proyecto.

\subsubsection{Vacío en la oferta existente}

El análisis de la oferta regional identifica dos actores relevantes en La Araucanía.

\begin{longtable}{@{}L{2.9cm}L{5.5cm}L{5.3cm}@{}}
\filaencab \encab{Proveedor} & \encab{Oferta} & \encab{Qué queda sin cubrir} \\ \midrule \endfirsthead
\filaencab \encab{Proveedor} & \encab{Oferta} & \encab{Qué queda sin cubrir} \\ \midrule \endhead
\textbf{IoT3.cl} &
Monitoreo y automatización en tiempo real: sensores, controladores y plataforma web. Fuerte presencia en
agricultura. Solución llave en mano por proyecto. &
Si el hardware ya está instalado y es de otra marca, no existe vía de integración. \\
\textbf{ControlTec.cl} (SpA, 2021) &
Telemetría industrial inalámbrica: niveles de estanques y combustible, climatización, agricultura.
Supervisión desde dispositivo móvil. &
Panel de telemetría acoplado a su propio hardware; no constituye una plataforma de ingesta abierta. \\
\bottomrule
\caption{Oferta regional identificada y espacio no cubierto.}
\label{tab:competencia}
\end{longtable}

\begin{cajaclave}[title={Espacio que este trabajo aborda}]
Ambos proveedores comercializan la solución completa acoplada a su propio hardware. Ninguno ofrece la
\textbf{capa de gestión desacoplada del equipamiento}: una plataforma neutral capaz de recibir datos de
sensores que la organización ya posee, cualquiera sea su procedencia.
\end{cajaclave}

\subsubsection{Relevancia académica}

La interoperabilidad entre protocolos IoT constituye un problema abierto: la proliferación de estándares
de aplicación no convergentes obliga a resolver la normalización en la capa de plataforma. El aporte de
este trabajo consiste en especificar y validar un mecanismo concreto de normalización \textbf{basado en
un estándar de industria ya existente y ampliamente desplegado}, en lugar de proponer un formato nuevo
que agregaría un elemento más a la fragmentación que pretende resolver.

\subsection{Partes interesadas}
\label{sec:interesados}

\begin{longtable}{@{}L{3.4cm}L{2.8cm}L{4.6cm}L{2.9cm}@{}}
\filaencab \encab{Parte interesada} & \encab{Rol} & \encab{Interés} & \encab{Levantamiento} \\ \midrule \endfirsthead
\filaencab \encab{Parte interesada} & \encab{Rol} & \encab{Interés} & \encab{Levantamiento} \\ \midrule \endhead
\datoUniversidad & Propietaria del proyecto & Enviar los datos de sus sensores a una interfaz de programación y administrarlos & \pendiente{Entrevista pendiente} \\
Responsable técnico de los sensores & Usuario administrador & Conocer qué dispositivos existen, si reportan y cómo incorporar uno nuevo & \pendiente{Por identificar} \\
Usuario operativo & Consumidor de información & Consultar mediciones y recibir avisos & \pendiente{Por identificar} \\
Pymes productivas de la región & Beneficiario potencial & Aprovechar sensores ya adquiridos sin rehacer la inversión & Caracterización indirecta \\
Profesor guía & Supervisor académico & Rigor metodológico y enfoque de redes & Reunión del 17/08/2026 \\
\bottomrule
\caption{Partes interesadas y estado de su levantamiento.}
\label{tab:interesados}
\end{longtable}

\begin{cajahallazgo}[title={Brecha declarada en el levantamiento}]
La rúbrica de evaluación exige incorporar \enquote{la percepción de las partes interesadas para validar
la relevancia del problema}. \textbf{Ese levantamiento directo está pendiente.} Las entrevistas al
responsable técnico de los sensores de la organización destinataria constituyen la acción prioritaria
del período, y sus resultados se incorporarán a esta sección y al catálogo de requerimientos.

Se opta por declarar esta brecha de forma explícita antes que presentar como levantada una percepción
que aún no ha sido recogida.
\end{cajahallazgo}

\subsection{Delimitación del alcance}

\begin{minipage}[t]{0.48\textwidth}
\textbf{El trabajo abarca:}
\begin{itemize}
\item La plataforma de ingesta, normalización, almacenamiento y gestión de datos de sensores.
\item La especificación del protocolo de comunicación y de la interfaz de programación asociada.
\item Adaptadores para al menos dos transportes distintos.
\item El diseño de red: segmentación, seguridad del transporte y monitoreo del estado del enlace.
\item La validación sobre un caso de uso con sensores reales.
\end{itemize}
\end{minipage}\hfill
\begin{minipage}[t]{0.48\textwidth}
\textbf{El trabajo no abarca:}
\begin{itemize}
\item Fabricación o comercialización de sensores.
\item Firmware de grado comercial, más allá del nodo de prueba necesario para validar.
\item Analítica avanzada, aprendizaje automático o predicción.
\item Despliegue productivo a escala ni operación continua posterior a la entrega.
\item Certificación normativa del equipamiento.
\end{itemize}
\end{minipage}
